Autonómne vozidlá sú v súčasnosti jedným z najrýchlejšie sa rozvíjajúcich technologických odvetví. Ide o formu mobility, ktorá je bezpečnejšia, efektívnejšia a pohodlnejšia než konvenčné spôsoby dopravy riadené človekom. Zavedenie autonómnych vozidiel vedie k presunu pracovnej sily od rutinných manuálnych úkonov smerom k činnostiam zameraným na monitorovanie, manažovanie a kontrolu systémov. Taktiež umožňujú ľuďom s fyzickým postihnutím samostatne operovať automobil. V súčasnosti sú vozidlá s čiastočnou autonómiou dostupné pre každého, kto túži po luxuse, ktorý tieto autá poskytujú, a má na to potrebné financie.\par
Autonómne vozidlá sa rozdeľujú do šiestich stupňov autonómie, od 0 (žiadna autonómia) až po 5 (úplná autonómia bez nutnosti zásahu človeka). Na nahradenie funkcií človeka využívajú senzory na vnímanie okolia a umelú inteligenciu na rozhodovanie. Technologické spoločnosti sa v tejto oblasti snažia o čo najefektívnejšie prepojenie týchto dvoch systémov.\par
Najväčšími bariérami pri integrácii autonómnych vozidiel do bežnej spoločnosti sú ich vysoká cena a neschopnosť navigovať v komplikovaných hraničných situáciách. Súčasný vývoj sa snaží tieto bariéry aktívne prekonávať.\par
Cieľom bakalárskej práce je navrhnúť a implementovať riešenie autonómneho vozidla, ktoré je schopné detekovať prekážky a adekvátne na ne reagovať. Projekt realizujeme na platforme Arduino s využitím knižnice na počítačové videnie OpenCV. V rámci práce preskúmame aktuálne prístupy k autonómii a zvážime životaschopnosť súčasných trendov. Výstupom bude praktická implementácia modelu vozidla a jeho testovanie v rôznych situáciách.\par
Výber témy bakalárskej práce vychádza z nášho záujmu o automatizáciu a inovácie v oblasti automobilov. Daná téma spája mnohé populárne technológie, ako je umelá inteligencia a automatizácia logistiky. Zároveň nás zaujala práca s mikrokontrolérmi a elektronikou. Práve kombináciou týchto oblastí sme sa pre túto tému nadchli.
